\chapter{유일한 빠른 출력}
% full version: chapters/13_muscles.tex

여러분이 세상에서 빠르게 하는 모든 일, 곧 말 한마디, 눈길 한 번, 걸음 하나는 근육의 수축이다. 기계에는 느린 화학적 출력이라는 두 번째 출력도 있는데, 이것은 16장에서 다룬다. 여기서는 근육을 본다.

사슬의 마지막 고리는 \nt{ACET} 뉴런이다. 근육 섬유 하나하나는 이런 뉴런 딱 하나의 말을 듣고, 뉴런 하나는 수백에서 2000개 정도의 섬유 무리를 제어한다. 흥미롭게도 신경과 근육의 연결은 몸에서 가장 믿을 만한 연결이다. 뇌 안의 연결은 신호의 절반 넘게를 잃지만, 여기서는 딸깍 한 번마다 필요한 양의 몇 배나 되는 물질이 뿜어진다. 출력에서 오류가 나면 동작을 망치므로, 신뢰성을 넉넉하게 사 둔 것이다.

\section*{북소리가 힘이 된다}

딸깍 한 번은 근육의 짧은 수축 한 번이다. 딸깍이 드물게 오면 근육은 따로따로 움찔한다. 자주 오면 수축들이 합쳐져 고른 힘이 된다. 북을 빠르게 두드리면 소리가 이어져 웅웅거리는 것과 같다. 근육은 가장 단순한 디지털-아날로그 변환기처럼 스파이크의 흐름을 힘으로 바꾼다.

\pencilfig{113}{%
% twitch = alpha function; the sum of twitches for spikes 1 s and 0.16 s apart (arbitrary units of time)
\begin{scope}[x=0.95cm, y=0.36cm]
\foreach \dx/\t in {0/드물게,4.6/자주}{
  \begin{scope}[xshift=\dx cm]
  \draw[pen] (0,0) -- (4,0); \node[font=\small\itshape, text=pcGray, anchor=south] at (2,5.4) {\t};
  \end{scope}}
\draw[penlite=pcRed] plot[smooth] coordinates {(0.00,0.00) (0.05,0.00) (0.10,0.00) (0.15,0.00) (0.20,0.00) (0.25,0.45) (0.30,0.73) (0.35,0.90) (0.40,0.98) (0.45,1.00) (0.50,0.98) (0.55,0.94) (0.60,0.88) (0.65,0.81) (0.70,0.74) (0.75,0.66) (0.80,0.59) (0.85,0.52) (0.90,0.46) (0.95,0.41) (1.00,0.35) (1.05,0.31) (1.10,0.27) (1.15,0.23) (1.20,0.20) (1.25,0.62) (1.30,0.88) (1.35,1.02) (1.40,1.08) (1.45,1.09) (1.50,1.06) (1.55,1.00) (1.60,0.93) (1.65,0.86) (1.70,0.78) (1.75,0.70) (1.80,0.62) (1.85,0.55) (1.90,0.48) (1.95,0.42) (2.00,0.37) (2.05,0.32) (2.10,0.28) (2.15,0.24) (2.20,0.21) (2.25,0.62) (2.30,0.88) (2.35,1.03) (2.40,1.09) (2.45,1.09) (2.50,1.06) (2.55,1.01) (2.60,0.94) (2.65,0.86) (2.70,0.78) (2.75,0.70) (2.80,0.62) (2.85,0.55) (2.90,0.48) (2.95,0.42) (3.00,0.37) (3.05,0.32) (3.10,0.28) (3.15,0.24) (3.20,0.21) (3.25,0.62) (3.30,0.88) (3.35,1.03) (3.40,1.09) (3.45,1.09) (3.50,1.06) (3.55,1.01) (3.60,0.94) (3.65,0.86) (3.70,0.78) (3.75,0.70) (3.80,0.62) (3.85,0.55) (3.90,0.48) (3.95,0.42) (4.00,0.37)};
\foreach \s in {0.20,1.20,2.20,3.20}{\draw[pcBlue, line width=0.8pt] (\s,-0.35) -- (\s,-1.2);}
\begin{scope}[xshift=4.6cm]
\draw[penlite=pcRed] plot[smooth] coordinates {(0.00,0.00) (0.05,0.00) (0.10,0.00) (0.15,0.00) (0.20,0.00) (0.25,0.45) (0.30,0.73) (0.35,0.90) (0.40,1.35) (0.45,1.68) (0.50,1.85) (0.55,2.19) (0.60,2.51) (0.65,2.64) (0.70,2.84) (0.75,3.13) (0.80,3.22) (0.85,3.27) (0.90,3.55) (0.95,3.62) (1.00,3.54) (1.05,3.82) (1.10,3.88) (1.15,3.80) (1.20,3.98) (1.25,4.06) (1.30,3.97) (1.35,4.07) (1.40,4.16) (1.45,4.09) (1.50,4.10) (1.55,4.23) (1.60,4.17) (1.65,4.10) (1.70,4.26) (1.75,4.23) (1.80,4.07) (1.85,4.27) (1.90,4.27) (1.95,4.12) (2.00,4.26) (2.05,4.29) (2.10,4.17) (2.15,4.24) (2.20,4.31) (2.25,4.21) (2.30,4.21) (2.35,4.31) (2.40,4.24) (2.45,4.16) (2.50,4.31) (2.55,4.27) (2.60,4.10) (2.65,4.30) (2.70,4.29) (2.75,4.15) (2.80,4.28) (2.85,4.31) (2.90,4.18) (2.95,4.25) (3.00,4.32) (3.05,4.22) (3.10,4.21) (3.15,4.32) (3.20,4.25) (3.25,4.17) (3.30,4.31) (3.35,4.27) (3.40,4.11) (3.45,3.86) (3.50,3.56) (3.55,3.25) (3.60,2.93) (3.65,2.63) (3.70,2.33) (3.75,2.06) (3.80,1.81) (3.85,1.58) (3.90,1.38) (3.95,1.19) (4.00,1.03)};
\foreach \s in {0.20,0.36,0.52,0.68,0.84,1.00,1.16,1.32,1.48,1.64,1.80,1.96,2.12,2.28,2.44,2.60,2.76,2.92,3.08,3.24}{\draw[pcBlue, line width=0.8pt] (\s,-0.35) -- (\s,-1.2);}
\end{scope}
\node[lbl, text=pcRed, anchor=west] at (1.2,1.9) {연축};
\node[lbl, text=pcRed, anchor=south] at (6.6,4.55) {고른 힘};
\node[lbl, text=pcBlue, anchor=north] at (4.3,-1.35) {뉴런의 딸깍};
\end{scope}
}{드문 딸깍은 근육의 연축을 하나씩 따로 일으키고, 잦은 딸깍은 합쳐져 고르고 훨씬 큰 힘이 된다.}

\section*{부동소수점 힘}

뇌에는 힘을 조절하는 다이얼이 둘 있다. 뉴런이 얼마나 자주 딸깍하는가, 그리고 섬유 무리가 몇 개나 켜지는가. 무리는 가장 약한 것부터 가장 강한 것까지 순서대로 켜지는데, 이 순서를 계산하는 주체는 없다. 작은 뉴런이 그저 더 쉽게 흥분할 뿐이다. 그 결과 새로 켜지는 무리는 저마다 이미 있는 힘에 거의 같은 \emph{비율}을 더한다. 달걀을 집을 때는 힘을 아주 작은 단위로 나누어 쓰고, 여행 가방을 들 때는 큰 단위로 쓴다. 엔지니어라면 여기서 부동소수점 수를 알아볼 것이다. 지수는 켜진 무리의 수가 정하고, 정확한 값은 딸깍의 빈도가 정한다. 그 대가로 동작이 강할수록 정밀도는 떨어진다. 영향력 있는 한 가설에 따르면 우리 동작이 그렇게 매끄러운 까닭이 이것이다. 매끄러운 명령이 흩어짐을 가장 작게 만든다.

\section*{당길 뿐 밀지 못한다}

근육은 당기기만 할 수 있다. 그래서 관절마다 서로 반대쪽으로 당기는 근육 한 쌍이 붙어 있다. 부호 대신 다시 전선 두 개를 쓰는 셈이다. 두 힘의 차이는 관절을 돌리고, 합은 관절을 단단하게 만든다. 가득 찬 찻잔을 나를 때는 두 근육에 한꺼번에 힘을 준다. 돌아가는 각도는 같지만 팔은 충격을 더 단단히 견딘다.

\pencilfig{13}{\pencilart{13}}{근육은 당기기만 하므로 두 개가 필요하다. 힘의 차이는 관절을 돌리고, 합은 관절을 단단하게 한다.}

\section*{늦게 오는 피드백}

근육에는 늘어남을 감지하는 감지기가 있고, 가장 짧은 루프는 척수 안에서 곧바로 닫힌다. 근육이 늘어나면 그 근육은 스스로 수축하고, 반대쪽 근육은 그동안 꺼진다. 반대쪽 근육을 끄는 것은 \nt{GLYC} 유형의 억제성 뉴런이다. 이 유형이 할 일이 여기 있었다.

그러나 모든 루프에는 지연이 있다. 척수에서는 수십 밀리초, 눈을 거치면 수백 밀리초다. 늦게 오는 피드백은 시스템을 흔든다. 도로를 늦게 보면서 핸들을 꺾는 운전자와 같다. 루프가 길수록 보정은 느려야 한다. 척수 루프가 흔들리기 시작하는 경계는 초당 8회 정도의 진동이다. 이 값은 건강한 사람이면 누구에게나 있는 미세한 떨림의 주파수와 가깝고, 신장 루프는 이 떨림의 유력한 원인 가운데 하나다.

그렇다면 우리는 어떻게 빠르고 정확하게 움직일까? 널리 퍼진 가설에 따르면 예측으로 해낸다. 몸의 내부 모델을 가지고, 감지기의 보고를 기다리지 않은 채 명령의 결과를 미리 계산한다.

\section*{메트로놈 없는 리듬}

걷기와 숨쉬기는 리듬이 있지만 메트로놈은 필요 없다. 서로를 억제하면서 차츰 지치는 두 뉴런 무리는 저절로 번갈아 일한다. 이긴 쪽이 지치면 쉬고 있던 상대가 우위를 차지하고, 이것이 되풀이된다. 위에서 내려오는 한결같은 명령 “걸어라”는 그 자리에서 “왼발, 오른발” 리듬으로 바뀐다.

\fullversion{13}
